%% kapa: exam=1 round=28 year=2017 period=2 form=A subject=appraisal_law subject_ko=감정평가관계법규 seq=2 number=2
%% source: Q-Net 감정평가사 기출문제 (cst00309) artlSeq=5199603 file=extracted/1008/5199603/2017년 제28회 감정평가사 1차시험 문제지/2교시 A형.pdf
%% conversion: kordoc 4.24.51 -> tools/gen1.py (auto); review=proofread-sonnet; flags=none
\begin{question}{2}
국토의 계획 및 이용에 관한 법령상 용도지구에 관한 설명으로 옳지 않은 것은?
\begin{choices}
  \item 용도지구는 토지의 이용 및 건축물의 용도 등에 대한 용도지역의 제한을 강화하거나 완화하여 적용함으로써 용도지역의 기능을 증진시키고 미관ㆍ경관ㆍ안전 등을 도모하기 위하여 지정한다.
  \item 용도지구의 지정 또는 변경은 도시ㆍ군관리계획으로 결정한다.
  \item 시설보호지구는 학교시설보호지구, 공용시설보호지구, 항만시설보호지구 및 공항시설보호지구로 세분된다.
  \item 개발제한구역안의 취락을 정비하기 위하여 필요한 지구는 자연취락지구이다.
  \item 주거기능, 공업기능, 유통․물류기능 및 관광․휴양기능 중 2 이상의 기능을 중심으로 개발․정비할 필요가 있는 경우 복합개발진흥지구로 지정한다.
\end{choices}
\end{question}
\answer{4}
